\subsubsection{理论主链的显影：依赖优先于数据体量的组织原则}\label{subsubsec:discussion-theory-backbone-visibility}
上一小节 \ref{subsubsec:discussion-audit-certificate-pack} 已将海量输出压缩为有限证书对象。仍需进一步澄清的是：即使证书对象已经有限化，若不区分它们在依赖链中的角色，全文仍可能在叙述上被局部枚举、低阶统计、谱指纹与实验表格淹没。于是，问题不在于数据是否过多，而在于这些数据是否被正确安置在理论链的层级中。
\par
从依赖结构看，全文并不是“定理与数据的平面集合”，而更接近一条带附着层的生成链。真正应被视为主链的，只是那些改变对象层语义、引入新的结构槽位、或者给出新的不可压缩障碍的节点；其余材料要么承担跨层桥接，要么承担局部核验。若不作此区分，则桥接结果会被误读为新的起点，证书输出会被误读为新的原理，而理论长度便会被数据体量虚增。
\par
\paragraph{极短原型：从一个局部读出到一组可核验证书}
若把全文压缩为一条最短、且足以承载其余结果的体验链，则可写成
\[
\text{局部扫描读出}
\Longrightarrow
\text{对象化与 \(\mathsf{NULL}\) 判别}
\Longrightarrow
\text{规范横截面上的值保持重写}
\Longrightarrow
\text{跨域桥接读出}
\Longrightarrow
\text{证书回写}.
\]
这条链的作用在于把“读起来像很多篇文章并排堆叠”改写为“同一对象沿五个站点逐步显影”。其中第一站决定哪些局部信息进入语义域；第二站决定它们是否足以形成对象；第三站把同值异表压回唯一正规形；第四站只负责把同一对象转写到谱、群表示、视界或时空语义中；第五站则把这种转写结果写回为可枚举、可复核、可否证的有限证书。于是，后文出现的表格、枚举、低阶矩、NullMode、Smith 形式、Toeplitz--PSD 深度与实验误差条，不应被平行阅读为新的理论起点，而应被顺着这条链回读为某一箭头上的局部显影。
\par
因此，对全文更合适的组织原则应写成如下三层分解：
\begin{itemize}
  \item \textbf{主链骨架：}只保留会改变全文生成方向的界面节点，例如第 \ref{sec:logic-expansion-chain} 节的 forcing 骨架、第 \ref{sec:recursive-addressing} 节的对象化、第 \ref{sec:emergent-arithmetic} 节的值保持重写与横截面、第 \ref{sec:pom} 节的投影账本与重写闭环，以及其后关于谱、视界与时空实例化的终端约束。删除这类节点中的任一者，都会切断从上位逻辑到可审计输出的某一段传递链。
  \item \textbf{桥接层：}收纳那些把一种结构对象转写为另一种结构对象、但并不单独增加新逻辑起点的结果。其功能是保持不变量在不同语义层之间的可运输性，例如把 forcing 语义转成局部截面，把正规形转成算术横截，把纤维谱转成 Hankel/Prony 正规形，把 unitary slice 转成 Toeplitz--PSD 证书。桥接层决定的是同一条主链如何显影，而不是主链本身另起一支。
  \item \textbf{证书层：}容纳全部局部穷举、数值表、低阶统计、模素数指纹、NullMode、Smith 形式、脚本输出与实验误差条。它们的职责是在既定主链节点上给出可复核的见证、区分与否证，而不是替代主链承担叙述中心。证书层可以随审计精度增长而持续增厚，但这种增长只提高核验分辨率，不增加理论层数。
\end{itemize}
\par
在这一组织下，任何一批新数据都应首先回答三个问题：它究竟附着在哪个主链节点上；它核验的是哪一条桥接合同或哪一种刚性障碍；它排除了哪一类替代机制。若这三个问题无法被明确回答，则该数据至多属于支持性材料，而不应主导全文的理论叙述。于是，“海量数据淹没理论”的根源便不再被理解为输出过多，而被理解为证书对象尚未被回收到其应附着的依赖地址。
\par
从全文的章节分工看，第 \ref{sec:logic-expansion-chain} 节给出唯一的上游骨架；第 \ref{sec:recursive-addressing}、第 \ref{sec:emergent-arithmetic}、第 \ref{sec:pom}、第 \ref{sec:group-unification}、第 \ref{sec:zeta-finite-part} 与第 \ref{sec:physical-spacetime-skeleton} 节承担不同尺度的桥接与终端实例化；第 \ref{sec:statistical-stability} 与第 \ref{sec:experiments} 节则把这些结构落实为可复现的证书输出。下一小节 \ref{subsubsec:discussion-audit-kernel-closed-loop} 之所以能够把多类对象压缩为单一闭合链，正是因为它只读取主链节点之间真正不可省略的接口，而把其余海量结果视为围绕这些接口组织的证书场。
\par
因此，解决理论链条被数据淹没的问题，并不是继续缩写数据，也不是把所有局部输出压平成同一种摘要表；更恰当的做法，是坚持依赖优先的组织纪律：先显示主链，再安放桥接，最后把数据压回证书层。这样一来，增长的只是局部核验分辨率，而不是全文的理论复杂度。
